% !PRINTED 269
\weberpage{269}
Es ist noch zu erw\"ahnen, da\ss{} die in §\,69, 4., 5. f\"ur die Invariantengleichung bei zusammengesetztem Transformationsgrade durchgef\"uhrte Betrachtung unver\"andert auch f\"ur die \gesperrt{Schlaefli}schen Modulargleichungen gilt, woraus wir schlie\ss{}en k\"onnen, da\ss{} alle diese Gleichungen rationale Zahlenkoeffizienten haben.

\begin{center}
\textbf{§\,75.\quad \gesperrt{Die irrationalen Formen der Modulargleichungen.}}
\end{center}

Den Transformationsgleichungen lassen sich durch Anwendung desselben Verfahrens weit einfachere Formen geben. Die Gleichungsformen, mit denen wir uns jetzt besch\"aftigen werden, enthalten die drei Funktionen $f$, $f_1$, $f_2$ zugleich; da man aber nach §\,34, (11), (12) zwei dieser Funktionen durch die dritte ausdr\"ucken kann, so lassen sich zwei von ihnen eliminieren, und man kann so zu den Gleichungen des vorigen Paragraphen gelangen. Wenn man f\"ur $f_1$, $f_2$ die Ausdr\"ucke durch $f$, oder f\"ur die drei Funktionen $f$, $f_1$, $f_2$ die Ausdr\"ucke durch $k^2$ einsetzt, so kommen Wurzelzeichen vor, woraus sich der Name dieser Gleichungen erkl\"art.

Wir setzen wie oben
\weq{1}{\begin{aligned}
u &= f(\omega), & v &= f\Bigl(\frac{c+\partial\,\omega}{a}\Bigr),\\
u_1 &= f_1(\omega), & v_1 &= \Bigl(\frac{2}{a}\Bigr) f_1\Bigl(\frac{c+\partial\,\omega}{a}\Bigr),\\
u_2 &= f_2(\omega), & v_2 &= \Bigl(\frac{2}{\partial}\Bigr) f_2\Bigl(\frac{c+\partial\,\omega}{a}\Bigr),\\
u\,u_1\,u_2 &= \sqrt{2}, & v\,v_1\,v_2 &= \Bigl(\frac{2}{n}\Bigr)\sqrt{2},
\end{aligned}}
und erhalten folgende zusammengeh\"orige Vertauschungen:
\weq{2}{\begin{array}{llll}
\omega, & uv, & u_1v_1, & u_2v_2,\\
-\dfrac{1}{\omega}, & \varrho\,uv, & \varrho\,u_2v_2, & \varrho\,u_1v_1,\\
\omega+1, & e^{-\frac{(n+1)\pi i}{24}}\sigma\,u_1v_1, & e^{-\frac{(n+1)\pi i}{24}}\sigma\,uv, & e^{\frac{(n+1)\pi i}{12}}\sigma\,u_2v_2,
\end{array}}
% Source scan and the row pattern confirm the exponent numerator (n+1).
worin $\varrho$, $\sigma$ die oben [§\,73, (4)] definierten dritten Einheitswurzeln sind.

Wir unterscheiden drei verschiedene F\"alle nach dem Verhalten von $n$ zum Modul 8.

1. $n+1\equiv 0$, (mod 8).
